\documentclass[11pt,a4paper]{article}
\usepackage{fontspec}
\setmainfont{Times New Roman}
\setsansfont{Arial}
\usepackage[russian]{babel}
\usepackage{amsmath,amssymb}
\usepackage[margin=19mm,headheight=15pt,headsep=7mm,footskip=10mm]{geometry}
\usepackage{xcolor,fancyhdr,titlesec,enumitem,hyperref}
\definecolor{muted}{HTML}{000000}
\hypersetup{unicode=true,colorlinks=true,linkcolor=black,pdftitle={Госэкзамен: 38 понятных доказательств}}
\pagestyle{fancy}\fancyhf{}
\fancyhead[L]{\small\sffamily\color{muted} Госэкзамен по математике · 38 доказательств}
\fancyfoot[L]{\footnotesize\color{muted} По темам программы · На основе конспекта 2021}
\fancyfoot[R]{\color{muted}\thepage}
\renewcommand{\headrulewidth}{0pt}
\setlength{\parindent}{0pt}\setlength{\parskip}{6pt}
\setlength{\abovedisplayskip}{7pt}\setlength{\belowdisplayskip}{7pt}
\setlist[enumerate]{leftmargin=19pt,itemsep=5pt,topsep=3pt,parsep=2pt}
\titleformat{\section}{\large\bfseries}{}{0pt}{}
\titlespacing*{\section}{0pt}{0pt}{10pt}
\emergencystretch=3em
\newcommand{\R}{\mathbb R}\newcommand{\C}{\mathbb C}
\newcommand{\E}{\mathbb E}\newcommand{\Var}{\operatorname{Var}}
\newcommand{\topic}[3]{\clearpage\section*{#1. #2}\phantomsection\addcontentsline{toc}{section}{#1. #2}{\small\color{muted}Конспект: #3.}\par}
\newcommand{\claim}{\textbf{Утверждение.} }
\newcommand{\idea}[1]{\par\textbf{Идея.} #1\par}
\newcommand{\uses}[1]{{\small\color{muted}\textbf{Используем:} #1.}\par}
\newcommand{\finish}{\hfill$\square$\par}
\newcommand{\remember}[1]{\par\textbf{Что запомнить.} #1}
\begin{document}
\pagecolor{white}\color{black}
\section*{По одному доказательству к каждой теме}
\textbf{38 разделов для подготовки к устному ответу.}

В каждом разделе выбрано одно важное утверждение из соответствующей темы. Условия выписаны перед доказательством; затем идут идея и последовательные шаги. Это подборка доказательств, а не замена всех определений и теорем каждого билета.

\textbf{Как учить.} Воспроизведите формулировку с условиями. Объясните идею одной фразой и восстановите шаги, не глядя в текст. Проверьте, где использованы условия.

Базовые результаты в строке «Используем» принимаются известными и не доказываются заново. В частности, раздел 34 опирается на теорему непрерывности Леви, а разделы 16 и 38 --- на интегральные теоремы.

Нумерация 1--38 соответствует программе госэкзамена 2025--2026. Ссылки относятся к напечатанным номерам страниц конспекта \textit{2021\_Oral (1) (2).pdf}. Доказательства переработаны: пропущенные переходы раскрыты, опечатки не перенесены.

В разделах 14 и 16 явно указаны удобные геометрические условия: прямоугольник и поверхность, заданная одним параметрическим куском. Это позволяет дать законченные доказательства без технического разбиения общей области.

\textbf{Обозначения.} $C^k$ --- функции с непрерывными производными до порядка $k$; $o(r)$ --- величина, отношение которой к $r$ стремится к нулю; $\|\cdot\|$ --- евклидова норма (для матриц --- согласованная операторная норма).
\clearpage
{\small\tableofcontents}

\topic{1}{Больцано--Вейерштрасс}{§ 1.1, с. 7--9}
\claim Из любой ограниченной вещественной последовательности $(x_n)$ можно выделить сходящуюся подпоследовательность.
\idea{Оставлять половину отрезка, содержащую бесконечно много членов. Длины отрезков стремятся к нулю.}
\uses{теорему о вложенных замкнутых отрезках}
\begin{enumerate}
\item Все $x_n$ лежат в некотором $I_0=[a,b]$. Если $a=b$, результат очевиден. Далее $a<b$.
\item Разделим $I_0$ пополам. Хотя бы одна половина содержит $x_n$ для бесконечно многих индексов: иначе их объединение содержало бы лишь конечное число индексов. Эту половину назовём $I_1$.
\item Повторяя выбор, получим замкнутые отрезки
\[
 I_0\supset I_1\supset I_2\supset\cdots,\qquad |I_k|=(b-a)2^{-k}.
\]
Существует единственная точка $c$, принадлежащая всем $I_k$.
\item Выбираем $n_k>n_{k-1}$ так, чтобы $x_{n_k}\in I_k$. Это возможно: подходящих индексов бесконечно много.
\item Обе точки $c$ и $x_{n_k}$ лежат в $I_k$, поэтому
\[
 |x_{n_k}-c|\le(b-a)2^{-k}\longrightarrow0.
\]
Следовательно, $x_{n_k}\to c$.
\end{enumerate}\finish
\remember{Индексы выбранных членов обязательно должны строго возрастать.}

\topic{2}{Достижение максимума на отрезке}{§ 1.2, с. 9--10}
\claim Непрерывная на $[a,b]$ вещественная функция ограничена и достигает максимума и минимума.
\idea{Выделить сходящуюся подпоследовательность точек и применить непрерывность.}
\uses{Больцано--Вейерштрасса; непрерывность; свойство супремума}
\begin{enumerate}
\item Если $f$ не ограничена, для каждого $n$ найдётся $x_n\in[a,b]$ с $|f(x_n)|>n$. Выделим $x_{n_k}\to c\in[a,b]$; принадлежность предела следует из замкнутости отрезка.
\item По непрерывности $f(x_{n_k})\to f(c)$, поэтому эта последовательность ограничена. Но $|f(x_{n_k})|>n_k\to\infty$. Противоречие: ограниченность доказана.
\item Существует конечное $M=\sup f([a,b])$. По свойству супремума выберем $y_n\in[a,b]$ так, что
\[
 M-1/n<f(y_n)\le M.
\]
Отсюда $f(y_n)\to M$.
\item Выделим $y_{n_k}\to d\in[a,b]$. Непрерывность даёт
\[
 f(d)=\lim_{k\to\infty}f(y_{n_k})=M.
\]
Максимум достигнут. Применяя доказанное к $-f$, получаем достижение минимума $f$.
\end{enumerate}\finish
\remember{Супремум --- число; для максимума нужна точка, в которой это число является значением функции.}

\topic{3}{Теорема о промежуточном значении}{§ 1.3, с. 11}
\claim Если $f\in C[a,b]$ и $f(a)<c<f(b)$, то существует $\xi\in(a,b)$ с $f(\xi)=c$. При обратном порядке значений вывод тот же.
\idea{Для $g=f-c$ всё время сохранять разные знаки на концах отрезка.}
\uses{теорему о вложенных отрезках; непрерывность}
\begin{enumerate}
\item Для $g(x)=f(x)-c$ имеем $g(a)<0<g(b)$.
\item Делим отрезок пополам. Если в середине $g=0$, всё доказано. Иначе выбираем половину, на левом конце которой $g<0$, а на правом $g>0$.
\item Повторяем, пока нуль не найден. Если процедура бесконечна, получаем $[a_n,b_n]$ с
\[
 g(a_n)<0<g(b_n),\qquad b_n-a_n=(b-a)2^{-n}.
\]
По теореме о вложенных отрезках $a_n\to\xi$ и $b_n\to\xi$ для одной точки $\xi$.
\item По непрерывности обе последовательности значений стремятся к $g(\xi)$. Предел отрицательных чисел не больше нуля, а положительных --- не меньше нуля. Значит, $g(\xi)\le0$ и $g(\xi)\ge0$, откуда $g(\xi)=0$.
\item Итак, $f(\xi)=c$. Поскольку $g(a)$ и $g(b)$ ненулевые, точка $\xi$ внутренняя. Обратный порядок знаков сводится к этому заменой $g$ на $-g$.
\end{enumerate}\finish
\remember{Без непрерывности функция может «перепрыгнуть» нужное значение.}

\topic{4}{Теорема Лагранжа}{§ 1.4.2, с. 12--13}
\claim Если $f\in C[a,b]$, $a<b$, и $f$ дифференцируема на $(a,b)$, то для некоторой $\xi\in(a,b)$
\[
 f(b)-f(a)=f'(\xi)(b-a).
\]
\idea{Вычесть прямую, соединяющую концы графика. У разности значения на концах равны.}
\uses{теоремы Вейерштрасса и Ферма}
\begin{enumerate}
\item Положим $k=(f(b)-f(a))/(b-a)$ и
\[
 h(x)=f(x)-f(a)-k(x-a).
\]
Тогда $h\in C[a,b]$, дифференцируема внутри и $h(a)=h(b)=0$.
\item Если $h\equiv0$, то $h'=0$ всюду внутри; можно выбрать любую внутреннюю точку.
\item Иначе $h$ принимает положительное или отрицательное значение. В первом случае её максимум положителен, во втором минимум отрицателен. Этот экстремум достигается по Вейерштрассу и не может лежать на концах, где $h=0$.
\item Значит, во внутренней точке $\xi$ по Ферма $h'(\xi)=0$. Так как $h'=f'-k$, имеем
\[
 f'(\xi)=\frac{f(b)-f(a)}{b-a}.
\]
\end{enumerate}\finish
\remember{Шаги 2--3 одновременно доказывают теорему Ролля для вспомогательной функции.}

\topic{5}{Тейлор с остатком Лагранжа}{§ 1.5.3, с. 15--17}
\claim Пусть $n\ge0$, $f\in C^n$ на отрезке между $a$ и $x\ne a$, а $f^{(n+1)}$ существует внутри. Для некоторой точки $\xi$ строго между $a$ и $x$
\[
 f(x)=\sum_{j=0}^n\frac{f^{(j)}(a)}{j!}(x-a)^j
 +\frac{f^{(n+1)}(\xi)}{(n+1)!}(x-a)^{n+1}.
\]
\idea{Подобрать последний коэффициент так, чтобы многократно применить Ролля.}
\uses{теорему Ролля}
\begin{enumerate}
\item Обозначим многочлен в сумме через $T_n$. Зафиксируем $x$ и положим
\[
 K=\frac{f(x)-T_n(x)}{(x-a)^{n+1}},\quad
 H(t)=f(t)-T_n(t)-K(t-a)^{n+1}.
\]
Имеем $H(x)=0$ и $H(a)=H'(a)=\cdots=H^{(n)}(a)=0$.
\item Из $H(a)=H(x)=0$ по Роллю получаем внутренний нуль $H'$. Если $n\ge1$, у $H'$ есть также нуль в $a$, поэтому между ними найдётся нуль $H''$.
\item Продолжая, после $n+1$ применений получаем $H^{(n+1)}(\xi)=0$ между $a$ и $x$. Гладкость в условии обеспечивает все применения Ролля. При $n=0$ достаточно первого.
\item Так как $T_n^{(n+1)}=0$, имеем
\[
 0=f^{(n+1)}(\xi)-K(n+1)!.
\]
Полученное $K=f^{(n+1)}(\xi)/(n+1)!$ подставляем в его определение.
\end{enumerate}\finish
\remember{$x$ зафиксирован; $K$ зависит от этого $x$, но не от переменной $t$.}

\topic{6}{Критерий неубывания через производную}{§ 1.6.1, с. 17}
\claim Для дифференцируемой на интервале $I$ функции равносильны: $f$ не убывает на $I$ и $f'(x)\ge0$ для всех $x\in I$. Если $f'>0$ всюду, функция строго возрастает.
\idea{В одну сторону --- знаки разностных отношений, в другую --- Лагранж для любых двух точек.}
\uses{определение производной; теорему Лагранжа}
\begin{enumerate}
\item \textbf{Необходимость.} Пусть $f$ не убывает. Фиксируем $x\in I$. При малом $h>0$ числитель $f(x+h)-f(x)\ge0$, поэтому
\[
 \frac{f(x+h)-f(x)}h\ge0.
\]
При $h<0$ числитель не больше нуля, знаменатель отрицателен: отношение снова неотрицательно. Переходя к пределу, получаем $f'(x)\ge0$.
\item \textbf{Достаточность.} Пусть $f'\ge0$. Для любых $u<v$ из $I$ функция непрерывна на $[u,v]$ по дифференцируемости. По Лагранжу для некоторой $\xi\in(u,v)$
\[
 f(v)-f(u)=f'(\xi)(v-u)\ge0.
\]
Это означает неубывание $f$.
\item Если всюду $f'>0$, последняя разность строго положительна: $f(u)<f(v)$ при любых $u<v$, то есть $f$ строго возрастает.
\end{enumerate}\finish
\remember{$f'>0$ --- достаточное, но не необходимое условие строгого возрастания: $x^3$ строго возрастает, хотя $f'(0)=0$.}

\topic{7}{Теорема Кантора}{§ 1.7, с. 19--20}
\claim Непрерывная $f:K\to\R$ на компакте $K\subset\R^m$ равномерно непрерывна:
\[
 \forall\varepsilon>0\ \exists\delta>0\ \forall x,y\in K:
 \|x-y\|<\delta\ \Longrightarrow\ |f(x)-f(y)|<\varepsilon.
\]
\idea{Если общей $\delta$ нет, найдутся близкие пары точек с далёкими значениями. Компактность приводит их к одной предельной точке.}
\uses{последовательную компактность; непрерывность}
\begin{enumerate}
\item Предположим противное. Тогда существует $\varepsilon_0>0$ такое, что для каждого $n$ можно выбрать $x_n,y_n\in K$ с
\[
 \|x_n-y_n\|<1/n,\qquad |f(x_n)-f(y_n)|\ge\varepsilon_0.
\]
\item Выделим $x_{n_k}\to a\in K$ по компактности.
\item По неравенству треугольника
\[
 \|y_{n_k}-a\|\le\|y_{n_k}-x_{n_k}\|+\|x_{n_k}-a\|\to0.
\]
Следовательно, $y_{n_k}\to a$ также.
\item По непрерывности обе последовательности значений стремятся к $f(a)$. Их разность стремится к нулю, что противоречит оценке снизу через $\varepsilon_0$.
\end{enumerate}\finish
\remember{Компактность обеспечивает предельную точку внутри $K$, где можно применить непрерывность.}

\topic{8}{Непрерывные частные производные}{§ 1.8, с. 20--22}
\claim Если все $f_{x_j}$ существуют около $a\in\R^m$ и непрерывны в $a$, то $f$ дифференцируема в $a$:
\[
 f(a+h)-f(a)=\sum_{j=1}^m f_{x_j}(a)h_j+o(\|h\|).
\]
\idea{Менять координаты по одной: каждое приращение станет одномерным.}
\uses{теорему Лагранжа; неравенство Коши--Буняковского}
\begin{enumerate}
\item Для малого $h$ положим $a^{(j)}=a+(h_1,\ldots,h_j,0,\ldots,0)$ и $a^{(0)}=a$. Тогда
\[
 f(a+h)-f(a)=\sum_{j=1}^m[f(a^{(j)})-f(a^{(j-1)})].
\]
\item В $j$-й разности меняется лишь $j$-я координата. По Лагранжу разность равна $f_{x_j}(c_j)h_j$, где $c_j$ лежит на соответствующем отрезке. При $h_j=0$ можно взять $c_j=a^{(j)}$.
\item Так как $\|c_j-a\|\le\|h\|$, по непрерывности $\alpha_j(h):=f_{x_j}(c_j)-f_{x_j}(a)\to0$.
\item Вычитая линейную часть, для остатка $r(h)$ получаем
\[
 |r(h)|=\left|\sum_j\alpha_j(h)h_j\right|
 \le\left(\sum_j|\alpha_j(h)|^2\right)^{1/2}\|h\|=o(\|h\|).
\]
Это требуемое определение дифференцируемости.
\end{enumerate}\finish
\remember{Одного существования частных производных недостаточно: в шаге 3 использована непрерывность в $a$.}

\topic{9}{Неявная функция одного уравнения}{§ 1.9, с. 22--23}
\claim Если $F\in C^1$ около $(a,b)\in\R^m\times\R$, $F(a,b)=0$, $F_y(a,b)\ne0$, то в некоторых окрестностях $W$ точки $a$ и $V$ точки $b$ уравнение $F(x,y)=0$ равносильно $y=\varphi(x)$, где $\varphi\in C^1(W)$ и
\[
 \varphi_{x_j}(x)=-F_{x_j}(x,\varphi(x))/F_y(x,\varphi(x)).
\]
\idea{К уравнению добавить координаты $x$ и обратить полученное отображение.}
\uses{теорему об обратной функции; цепное правило}
\begin{enumerate}
\item Для $\Phi(x,y)=(x,F(x,y))$ имеем
\[
 D\Phi=\begin{pmatrix}I_m&0\\F_x&F_y\end{pmatrix},\qquad
 \det D\Phi(a,b)=F_y(a,b)\ne0.
\]
По теореме об обратной функции $\Phi$ имеет локальное $C^1$-обратное $\Psi$ около $(a,0)$.
\item Первые $m$ координат $\Phi$ не меняет, поэтому $\Psi(u,v)=(u,g(u,v))$. Положим $\varphi(x)=g(x,0)$: тогда $F(x,\varphi(x))=0$ и $\varphi(a)=b$.
\item Выбираем $W\times V$ внутри области взаимной обратимости. По непрерывности, уменьшая $W$, обеспечиваем $(x,0)$ в области $\Psi$ и $\varphi(x)\in V$ для всех $x\in W$.
\item Если $F(x,y)=0$ в $W\times V$, то $\Phi(x,y)=(x,0)$; единственность обратного образа даёт $y=\varphi(x)$. Обратное уже получено в шаге 2.
\item Дифференцируя $F(x,\varphi(x))=0$, получаем $F_{x_j}+F_y\varphi_{x_j}=0$. По непрерывности $F_y$ не обращается в нуль в достаточно малой окрестности, и деление допустимо.
\end{enumerate}\finish
\remember{Здесь доказывается теорема о неявной функции, а теорема об обратной функции принимается известной.}

\topic{10}{Положительный гессиан даёт минимум}{§ 1.10, с. 23--25}
\claim Если $f\in C^2$ около $a\in\R^m$, $\nabla f(a)=0$, а $H=(f_{x_ix_j}(a))$ положительно определена, то $a$ --- точка строгого локального минимума.
\idea{Квадратичный член Тейлора положителен равномерно по направлениям, а остаток меньше него.}
\uses{Тейлора второго порядка; Вейерштрасса на единичной сфере}
\begin{enumerate}
\item Функция $u\mapsto u^THu$ непрерывна и положительна на единичной сфере. Сфера компактна, поэтому минимум достигается и равен некоторому $\mu>0$.
\item Записывая $h=\|h\|u$, где $\|u\|=1$, получаем
\[
 h^THh\ge\mu\|h\|^2.
\]
\item По Тейлору с учётом $\nabla f(a)=0$
\[
 f(a+h)-f(a)=\tfrac12h^THh+r(h),\qquad r(h)=o(\|h\|^2).
\]
По определению малого $o$ выберем $\delta>0$ так, чтобы при $0<\|h\|<\delta$ выполнялось $|r(h)|\le\mu\|h\|^2/4$.
\item Тогда
\[
 f(a+h)-f(a)\ge\frac\mu2\|h\|^2-\frac\mu4\|h\|^2
 =\frac\mu4\|h\|^2>0.
\]
Это определение строгого локального минимума.
\end{enumerate}\finish
\remember{Общую по всем направлениям константу $\mu>0$ даёт именно компактность сферы.}
\topic{11}{Производная интеграла с переменным пределом}{§ 1.11.2, с. 26--27}
\claim Пусть $f$ интегрируема по Риману на $[a,b]$ и непрерывна в $x_0\in(a,b)$. Тогда для $F(x)=\int_a^x f(t)\,dt$ существует $F'(x_0)=f(x_0)$.
\idea{Разностное отношение --- среднее значение $f$ на коротком отрезке около $x_0$. Непрерывность приближает это среднее к $f(x_0)$.}
\uses{аддитивность и оценку интеграла; определение непрерывности}
\begin{enumerate}
\item Для $h\ne0$, достаточно малого, аддитивность даёт
\[
 \frac{F(x_0+h)-F(x_0)}h-f(x_0)
 =\frac1h\int_{x_0}^{x_0+h}[f(t)-f(x_0)]\,dt.
\]
При отрицательном $h$ интеграл ориентирован: $\int_v^u=-\int_u^v$.
\item Зафиксируем $\varepsilon>0$. По непрерывности существует $\delta>0$ такое, что $|f(t)-f(x_0)|<\varepsilon$ при $|t-x_0|<\delta$.
\item Если $0<|h|<\delta$, эта оценка действует на всём отрезке между $x_0$ и $x_0+h$. Поэтому при любом знаке $h$
\[
 \left|\frac1h\int_{x_0}^{x_0+h}[f(t)-f(x_0)]\,dt\right|
 \le\frac1{|h|}\,\varepsilon|h|=\varepsilon.
\]
\item Значит, разностное отношение стремится к $f(x_0)$, что доказывает утверждение.
\end{enumerate}\finish
\remember{Непрерывность $f$ нужна в точке, где ищем производную. Интегрируемость на всём отрезке обеспечивает определённость $F$.}

\topic{12}{Непрерывность равномерного предела}{§ 1.12.3, с. 30--31}
\claim Если непрерывные на $E\subset\R^m$ функции $f_n$ сходятся к $f$ равномерно на $E$, то $f$ непрерывна на $E$ (в относительном смысле).
\idea{Вставить одну достаточно хорошую непрерывную функцию между двумя значениями предела. Получатся три малые разности.}
\uses{неравенство треугольника; определения равномерной сходимости и непрерывности}
\begin{enumerate}
\item Фиксируем $a\in E$ и $\varepsilon>0$. Равномерная сходимость позволяет выбрать один номер $N$ так, что
\[
 |f_N(x)-f(x)|<\varepsilon/3\quad\text{для всех }x\in E.
\]
В частности, это верно и в $a$.
\item Функция $f_N$ непрерывна в $a$. Для уже выбранного $N$ существует $\delta>0$ такое, что
\[
 x\in E,\quad\|x-a\|<\delta
 \quad\Longrightarrow\quad |f_N(x)-f_N(a)|<\varepsilon/3.
\]
\item Для таких $x$ применим треугольник:
\begin{align*}
 |f(x)-f(a)|&\le |f(x)-f_N(x)|+|f_N(x)-f_N(a)|\\
 &\quad+|f_N(a)-f(a)|<\varepsilon.
\end{align*}
Это непрерывность $f$ в $a$. Поскольку $a$ произвольно, утверждение доказано на всём $E$.
\end{enumerate}\finish
\remember{Порядок выбора: сначала $N$, подходящий для всех точек, затем $\delta$ для фиксированной функции $f_N$. При поточечной сходимости первый шаг не работает.}

\topic{13}{Сходимость степенного ряда внутри круга}{§ 1.13.1, с. 33--34}
\claim Если $\sum_{n=0}^\infty a_n(x-x_0)^n$ сходится в точке $x_1\ne x_0$, то для каждого $0<r<|x_1-x_0|$ он сходится абсолютно и равномерно при $|x-x_0|\le r$.
\idea{Сходимость в одной точке ограничивает члены; ближе к центру появляется геометрический множитель меньше единицы.}
\uses{ограниченность членов сходящегося ряда; геометрический ряд; признак Вейерштрасса}
\begin{enumerate}
\item Из сходимости следует $a_n(x_1-x_0)^n\to0$. Всякая сходящаяся последовательность ограничена, поэтому существует $M$ с
\[
 |a_n(x_1-x_0)^n|\le M\quad\text{для всех }n.
\]
\item Положим $q=r/|x_1-x_0|<1$. Для $|x-x_0|\le r$ получаем
\[
 |a_n(x-x_0)^n|
 =|a_n(x_1-x_0)^n|\left|\frac{x-x_0}{x_1-x_0}\right|^n
 \le Mq^n.
\]
\item Числовой ряд $\sum Mq^n=M/(1-q)$ сходится и не зависит от $x$. По признаку Вейерштрасса исходный ряд сходится абсолютно и равномерно на указанном отрезке.
\end{enumerate}\finish
\remember{Утверждение ничего не обещает в граничных точках $|x-x_0|=|x_1-x_0|$: там $q=1$ и геометрическая оценка перестаёт работать. Это основа понятия радиуса сходимости.}

\topic{14}{Критерий потенциальности на прямоугольнике}{§ 1.16, с. 43--44}
\claim Пусть $D=(a,b)\times(c,d)$ и $P,Q\in C^1(D)$. Поле $(P,Q)$ имеет потенциал $\varphi$, то есть $\varphi_x=P$, $\varphi_y=Q$, тогда и только тогда, когда $P_y=Q_x$ в $D$.
\idea{Построить потенциал как интеграл по ломаной: сначала по горизонтали, затем по вертикали.}
\uses{дифференцирование интеграла; равенство смешанных производных}
\begin{enumerate}
\item \textbf{Необходимость.} Если $\varphi_x=P$ и $\varphi_y=Q$, то из $P,Q\in C^1$ следует $\varphi\in C^2$. Поэтому
\[
 P_y=\varphi_{xy}=\varphi_{yx}=Q_x.
\]
\item \textbf{Достаточность.} Пусть $P_y=Q_x$. Фиксируем $(x_0,y_0)\in D$ и положим
\[
 \varphi(x,y)=\int_{x_0}^x P(s,y_0)\,ds+\int_{y_0}^y Q(x,t)\,dt.
\]
Оба отрезка лежат в $D$: здесь используется форма прямоугольника.
\item По производной интеграла $\varphi_y(x,y)=Q(x,y)$.
\item Дифференцируя по $x$ и используя $Q_x=P_y$, получаем
\begin{align*}
 \varphi_x(x,y)&=P(x,y_0)+\int_{y_0}^y Q_x(x,t)\,dt\\
 &=P(x,y_0)+\int_{y_0}^y P_y(x,t)\,dt=P(x,y).
\end{align*}
Значит, построенная $\varphi$ является потенциалом.
\end{enumerate}\finish
\remember{Здесь доказан критерий на прямоугольнике. На произвольной области равенство $P_y=Q_x$ само по себе не гарантирует глобального потенциала; в общем критерии нужна односвязность.}

\topic{15}{Дивергенция и нулевой поток}{§ 1.17, с. 44}
\claim Для поля $F\in C^1(G,\R^3)$ в открытой области $G$ равносильны: $\operatorname{div}F=0$ всюду и поток через границу любого замкнутого шара $\overline B\subset G$ равен нулю. При $\operatorname{div}F=0$ нулевой поток также имеет любая допустимая замкнутая поверхность, ограничивающая объём в $G$.
\idea{Гаусс превращает поток в интеграл дивергенции. Ненулевое непрерывное значение сохраняет знак в маленьком шаре.}
\uses{формулу Остроградского--Гаусса; непрерывность дивергенции}
\begin{enumerate}
\item Если $\operatorname{div}F=0$, по Гауссу для любого объёма $V$ с кусочно-гладкой границей и $\overline V\subset G$
\[
 \iint_{\partial V}(F,n)\,dS=\iiint_V\operatorname{div}F\,dV=0.
\]
Нормаль здесь внешняя. В частности, утверждение верно для каждого шара.
\item Обратно, пусть все потоки через сферы равны нулю. Допустим, в некоторой $a\in G$ имеем $d:=\operatorname{div}F(a)>0$.
\item По непрерывности найдётся шар $\overline B(a,r)\subset G$, где $\operatorname{div}F>d/2$. Тогда
\[
 \iint_{\partial B}(F,n)\,dS
 =\iiint_B\operatorname{div}F\,dV
 \ge(d/2)\operatorname{Vol}(B)>0,
\]
что противоречит условию. Если $d<0$, тот же довод применяем к $-F$.
\item Следовательно, нигде нет ненулевой дивергенции: $\operatorname{div}F=0$.
\end{enumerate}\finish
\remember{Не доказываем Гаусса заново: доказываем важное следствие, объясняющее смысл соленоидальности. Одной выбранной поверхности для обратного вывода недостаточно.}

\topic{16}{Стокс для параметрической поверхности}{§ 1.15, с. 42--43}
\claim Пусть $r(u,v)$ --- регулярная $C^2$-параметризация поверхности $S$ над замыканием ограниченной плоской области $D$ с кусочно-гладкой границей, взаимно однозначная на $\overline D$. Для $F\in C^1$ около $S$, ориентации $r_u\times r_v$ и согласованной границы
\[
 \oint_{\partial S}(F,dr)=\iint_S(\operatorname{rot}F,n)\,dS.
\]
\idea{Перенести криволинейный интеграл на плоскость параметров и применить Грина.}
\uses{формулу Грина; цепное правило; свойства смешанного произведения}
\begin{enumerate}
\item При $dr=r_u\,du+r_v\,dv$ положим $A=F(r)\cdot r_u$, $B=F(r)\cdot r_v$. Тогда по определению интеграла
\[
 \oint_{\partial S}F\cdot dr=\oint_{\partial D}A\,du+B\,dv
 =\iint_D(B_u-A_v)\,du\,dv.
\]
Граница $\partial D$ ориентирована положительно.
\item Цепное правило даёт
\[
 B_u-A_v=(DF(r)r_u)\cdot r_v-(DF(r)r_v)\cdot r_u.
\]
Члены $F(r)\cdot r_{vu}$ и $F(r)\cdot r_{uv}$ сократились, поскольку $r\in C^2$.
\item Для любых векторов $p,q$ верно тождество
\[
 (DF\,p)\cdot q-(DF\,q)\cdot p
 =(\operatorname{rot}F)\cdot(p\times q).
\]
Оно проверяется раскрытием трёх координат. С $p=r_u$, $q=r_v$ получаем нужный подынтегральный член.
\item По определению ориентированного поверхностного интеграла $n\,dS=(r_u\times r_v)\,du\,dv$. Подстановка завершает доказательство.
\end{enumerate}\finish
\remember{Доказана формула для одного параметрического куска. Для поверхности из нескольких согласованных кусков внутренние граничные интегралы при сложении сокращаются.}

\topic{17}{Признак Дини сходимости ряда Фурье}{§ 1.18.2--1.18.3, с. 45--47}
\claim Пусть $f$ --- $2\pi$-периодическая, $f\in L^1[-\pi,\pi]$. Если для некоторых $S$ и $\delta>0$
\[
 \int_0^\delta\frac{|f(x+t)+f(x-t)-2S|}{t}\,dt<\infty,
\]
то частичные суммы её ряда Фурье удовлетворяют $S_n(f,x)\to S$.
\idea{Представить ошибку как осциллирующий интеграл от одной интегрируемой функции.}
\uses{ядро Дирихле; лемму Римана--Лебега: для $g\in L^1(0,\pi)$ интеграл $\int_0^\pi g(t)\sin(\lambda t)\,dt\to0$ при $\lambda\to\infty$}
\begin{enumerate}
\item В нормировке конспекта
\[
 D_n(t)=\tfrac12+\sum_{k=1}^n\cos kt
 =\frac{\sin((n+1/2)t)}{2\sin(t/2)},\qquad
 \int_0^\pi D_n(t)\,dt=\frac\pi2.
\]
Формула свёртки для частичной суммы даёт
\[
 S_n(f,x)-S=\frac1\pi\int_0^\pi
 g(t)\sin((n+1/2)t)\,dt,
\]
где $g(t)=[f(x+t)+f(x-t)-2S]/[2\sin(t/2)]$.
\item При $0<t\le\pi$ имеем $\sin(t/2)\ge t/\pi$. Поэтому около нуля
\[
 |g(t)|\le\frac\pi2\frac{|f(x+t)+f(x-t)-2S|}{t}.
\]
Правая часть интегрируема по условию (при необходимости уменьшаем $\delta$ до числа меньше $\pi$).
\item Вдали от нуля знаменатель отделён от нуля; числитель интегрируем по периодичности и $f\in L^1$. Итак, $g\in L^1(0,\pi)$.
\item По лемме Римана--Лебега интеграл ошибки стремится к нулю. Значит, $S_n(f,x)\to S$.
\end{enumerate}\finish
\remember{Основная работа --- проверить интегрируемость $g$ около нуля; именно это обеспечивает условие Дини.}

\topic{18}{Равномерная сходимость для гладкой периодической функции}{§ 1.19, с. 47--48}
\claim Если $f$ --- вещественная $2\pi$-периодическая функция класса $C^1(\R)$, её ряд Фурье сходится к $f$ абсолютно и равномерно на $\R$.
\idea{Производная даёт множитель $1/n$ в коэффициентах. Неравенства Бесселя и Коши--Буняковского делают их сумму конечной.}
\uses{интегрирование по частям; неравенство Бесселя; признак Вейерштрасса; раздел 17}
\begin{enumerate}
\item Пусть $a_n,b_n$ --- коэффициенты Фурье $f$ в нормировке $1/\pi$, а $A_n,B_n$ --- коэффициенты $f'$. Интегрирование по частям и $f(-\pi)=f(\pi)$ дают
\[
 A_n=nb_n,\qquad B_n=-na_n\quad(n\ge1).
\]
Граничные члены сокращаются: синусы на концах нулевые, косинусы равны.
\item По Бесселю для $f'\in L^2[-\pi,\pi]$
\[
 \sum_{n=1}^\infty n^2(a_n^2+b_n^2)
 \le\frac1\pi\int_{-\pi}^\pi|f'(x)|^2\,dx<\infty.
\]
\item Так как $(|a_n|+|b_n|)^2\le2(a_n^2+b_n^2)$, по Коши--Буняковскому
\[
 \sum_{n=1}^\infty(|a_n|+|b_n|)
 \le\left(2\sum_{n=1}^\infty n^2(a_n^2+b_n^2)\right)^{1/2}
 \left(\sum_{n=1}^\infty\frac1{n^2}\right)^{1/2}<\infty.
\]
Отсюда по Вейерштрассу ряд Фурье абсолютно и равномерно сходится к некоторой функции $g$.
\item Производная $f'$ ограничена по периодичности. По Лагранжу $|f(x\pm t)-f(x)|\le M|t|$, поэтому условие Дини с $S=f(x)$ выполнено в каждой точке. По разделу 17 тот же ряд сходится к $f(x)$. Единственность предела даёт $g=f$.
\end{enumerate}\finish
\remember{Условие $C^1$ сильнее общих условий Гёльдера, зато даёт короткое доказательство даже абсолютной сходимости.}

\topic{19}{Непрерывность преобразования Фурье}{§ 1.20.2, с. 49--51}
\claim Если $f\in L^1(\R)$, то
\[
 \widehat f(y)=\frac1{\sqrt{2\pi}}\int_{\R}f(x)e^{-ixy}\,dx
\]
равномерно непрерывна на $\R$.
\idea{На хвосте мал интеграл $|f|$, а на ограниченном отрезке мало приращение экспоненты. Оценка не зависит от $y$.}
\uses{оценку $|e^{it}-1|\le|t|$; малость хвоста абсолютно сходящегося интеграла}
\begin{enumerate}
\item Обозначим $C=\int_\R|f|$. Из $|e^{-ixy}|=1$ следует существование интеграла. Для любых $y,h$
\[
 \sqrt{2\pi}\,|\widehat f(y+h)-\widehat f(y)|
 \le\int_\R|f(x)|\,|e^{-ixh}-1|\,dx.
\]
\item Пусть $\varepsilon>0$. Выберем $R>0$ так, чтобы
\[
 2\int_{|x|>R}|f(x)|\,dx<\tfrac12\varepsilon\sqrt{2\pi}.
\]
Это возможно по абсолютной интегрируемости. При $|x|>R$ используем $|e^{-ixh}-1|\le2$.
\item При $|x|\le R$ имеем $|e^{-ixh}-1|\le R|h|$, и соответствующая часть интеграла не превосходит $RC|h|$.
\item Если $C=0$, преобразование тождественно нулевое. Иначе выбираем $|h|<\varepsilon\sqrt{2\pi}/(2RC)$. Обе части вместе меньше $\varepsilon\sqrt{2\pi}$, значит, разность преобразований меньше $\varepsilon$.
\end{enumerate}\finish
\remember{Выбранный предел на $|h|$ не зависит от $y$: доказана равномерная, а не только поточечная непрерывность.}

\topic{20}{Расстояние от точки до плоскости}{§ 2.1.3, с. 53--54}
\claim Для плоскости $\Pi:Ax+By+Cz+D=0$, где $n=(A,B,C)\ne0$, и точки $M=(x_0,y_0,z_0)$
\[
 d(M,\Pi)=\frac{|Ax_0+By_0+Cz_0+D|}{\sqrt{A^2+B^2+C^2}}.
\]
\idea{Найти основание перпендикуляра вдоль нормали и доказать, что оно ближе всех остальных точек плоскости.}
\uses{скалярное произведение; теорему Пифагора}
\begin{enumerate}
\item Ищем $H=M-tn$. Условие $H\in\Pi$ равносильно
\[
 n\cdot M-t\|n\|^2+D=0,
 \quad t=\frac{n\cdot M+D}{\|n\|^2}.
\]
Деление допустимо, поскольку $n\ne0$.
\item Вектор $M-H=tn$ перпендикулярен плоскости. Для любой $X\in\Pi$ имеем $n\cdot(X-H)=0$, так как обе точки удовлетворяют уравнению плоскости.
\item Поэтому $M-H$ и $H-X$ ортогональны, и
\[
 \|M-X\|^2=\|M-H\|^2+\|H-X\|^2\ge\|M-H\|^2.
\]
Минимум расстояния достигается именно при $X=H$.
\item Наконец,
\[
 d(M,\Pi)=\|M-H\|=|t|\|n\|
 =\frac{|n\cdot M+D|}{\|n\|},
\]
что и даёт формулу в координатах.
\end{enumerate}\finish
\remember{Недостаточно найти длину нормальной проекции: шаг 3 объясняет, почему это действительно наименьшее расстояние.}
\topic{21}{Фокальное свойство эллипса}{§ 2.2.2, с. 56--57}
\claim Для $a>b>0$ и $c=\sqrt{a^2-b^2}$ каждая точка эллипса $x^2/a^2+y^2/b^2=1$ имеет сумму расстояний до фокусов $F_1=(-c,0)$, $F_2=(c,0)$, равную $2a$.
\idea{Подставить уравнение эллипса в квадраты расстояний: получатся полные квадраты.}
\uses{формулу расстояния между точками; уравнение эллипса}
\begin{enumerate}
\item Пусть $M=(x,y)$ лежит на эллипсе. Тогда $|x|\le a$ и $y^2=b^2(1-x^2/a^2)$. Вычислим $r_1=|MF_1|$:
\begin{align*}
 r_1^2&=(x+c)^2+y^2\\
 &=x^2+2cx+c^2+b^2-\frac{b^2}{a^2}x^2\\
 &=\frac{c^2}{a^2}x^2+2cx+a^2
 =\left(a+\frac ca x\right)^2.
\end{align*}
Здесь использованы $c^2+b^2=a^2$ и $1-b^2/a^2=c^2/a^2$.
\item Аналогично
\[
 r_2^2=(x-c)^2+y^2=\left(a-\frac ca x\right)^2.
\]
\item При извлечении корня нужно проверить знаки. Так как $c<a$ и $|x|\le a$,
\[
 a\pm\frac ca x\ge a-c>0.
\]
Следовательно, $r_1=a+cx/a$ и $r_2=a-cx/a$.
\item Складывая, получаем $r_1+r_2=2a$.
\end{enumerate}\finish
\remember{Доказано фокальное свойство точек канонического эллипса. Проверка знаков перед извлечением корня обязательна: вообще $\sqrt{u^2}=|u|$.}

\topic{22}{Теорема Кронекера--Капелли}{§ 2.3.2, с. 57--58}
\claim Система $Ax=b$, где $A$ --- матрица размера $m\times n$, совместна тогда и только тогда, когда $\operatorname{rank}A=\operatorname{rank}(A\mid b)$.
\idea{Решение существует ровно тогда, когда добавленный столбец уже выражается через старые и не увеличивает их линейную оболочку.}
\uses{ранг как размерность пространства столбцов; свойство размерности вложенных подпространств}
\begin{enumerate}
\item Обозначим столбцы $A$ через $a_1,\ldots,a_n\in\R^m$. Матричное уравнение $Ax=b$ означает
\[
 x_1a_1+\cdots+x_na_n=b.
\]
Следовательно, система совместна ровно при $b\in L$, где $L=\operatorname{span}(a_1,\ldots,a_n)$.
\item Пространство столбцов расширенной матрицы равно $L'=\operatorname{span}(L,b)$; всегда $L\subset L'$.
\item Если решение существует, то $b\in L$. Поэтому $L'=L$ и ранги совпадают.
\item Обратно, равенство рангов означает $\dim L=\dim L'$. Для конечномерных пространств включение и равенство размерностей дают $L=L'$: иначе добавление вектора из $L'\setminus L$ к базису $L$ увеличило бы размерность.
\item Значит, $b\in L$. По определению линейной оболочки существуют коэффициенты $x_1,\ldots,x_n$, выражающие $b$ через столбцы $A$. Это и есть решение системы.
\end{enumerate}\finish
\remember{Равенство рангов даёт существование решения. Единственность --- отдельный вопрос: для неё нужен ещё полный столбцовый ранг $n$.}

\topic{23}{Теорема о размерности ядра и образа}{§ 2.4.2, с. 59--60}
\claim Для линейного $T:U\to V$ с конечномерным $U$
\[
 \dim U=\dim\ker T+\dim\operatorname{Im}T.
\]
\idea{Дополнить базис ядра до базиса всего пространства; образы добавленных векторов образуют базис образа.}
\uses{дополнение линейно независимой системы до базиса}
\begin{enumerate}
\item Пусть $e_1,\ldots,e_k$ --- базис $\ker T$. Дополним его векторами $e_{k+1},\ldots,e_n$ до базиса $U$. Здесь $n=\dim U$.
\item Любой $u\in U$ имеет вид $\sum_{j=1}^n\alpha_je_j$. Так как $Te_j=0$ при $j\le k$, имеем
\[
 Tu=\sum_{j=k+1}^n\alpha_jTe_j.
\]
Значит, $Te_{k+1},\ldots,Te_n$ порождают образ.
\item Докажем их независимость. Если $\sum_{j=k+1}^n\beta_jTe_j=0$, то $w=\sum_{j=k+1}^n\beta_je_j\in\ker T$. Поэтому также $w=\sum_{j=1}^k\gamma_je_j$.
\item Вычитая два выражения $w$, получаем линейную комбинацию полного базиса, равную нулю. По независимости базиса все $\beta_j$ равны нулю. Независимость образов доказана.
\item Таким образом, образ имеет базис из $n-k$ векторов. Следовательно, $\dim\operatorname{Im}T=n-k$, откуда формула.
\end{enumerate}\finish
\remember{Самая важная часть --- независимость образов. Она неверна для произвольно выбранных векторов: здесь помогает дополнение именно базиса ядра.}

\topic{24}{Независимость собственных векторов}{§ 2.5, с. 60--61}
\claim Собственные векторы $v_1,\ldots,v_k$ линейного оператора $T$, отвечающие попарно различным собственным значениям $\lambda_1,\ldots,\lambda_k$, линейно независимы.
\idea{Применить $T-\lambda_k I$: последний вектор исчезнет, остальные сохранятся с ненулевыми коэффициентами.}
\uses{математическую индукцию; $Tv_j=\lambda_jv_j$, $v_j\ne0$}
\begin{enumerate}
\item При $k=1$ утверждение верно, поскольку собственный вектор ненулевой.
\item Предположим утверждение для $k-1$ векторов и рассмотрим равенство
\[
 \alpha_1v_1+\cdots+\alpha_kv_k=0.
\]
Применяем линейный оператор $T-\lambda_k I$ и получаем
\[
 \sum_{j=1}^{k-1}\alpha_j(\lambda_j-\lambda_k)v_j=0.
\]
\item По предположению индукции первые $k-1$ векторов независимы, поэтому $\alpha_j(\lambda_j-\lambda_k)=0$ для всех $j<k$. Значения различны, значит, $\alpha_j=0$ при $j<k$.
\item Исходное равенство превращается в $\alpha_kv_k=0$. Поскольку $v_k\ne0$, получаем $\alpha_k=0$.
\end{enumerate}\finish
\remember{Следствие: в $n$-мерном пространстве $n$ различных собственных значений дают базис собственных векторов и диагонализируемость. Обратное требование не обязательно: у диагонализируемого оператора значения могут повторяться.}

\topic{25}{Ортогональность собственных векторов}{§ 2.6, с. 62--63}
\claim Если $T$ самосопряжён в вещественном евклидовом пространстве, то собственные векторы, отвечающие различным собственным значениям, ортогональны.
\idea{Вычислить одно скалярное произведение двумя способами, перенося оператор между аргументами.}
\uses{определение самосопряжённости: $(Tx,y)=(x,Ty)$}
\begin{enumerate}
\item Пусть $Tv=\lambda v$, $Tw=\mu w$, где $\lambda\ne\mu$. Тогда по линейности скалярного произведения
\[
 (Tv,w)=\lambda(v,w),\qquad (v,Tw)=\mu(v,w).
\]
\item Самосопряжённость означает равенство левых частей. Следовательно,
\[
 (\lambda-\mu)(v,w)=0.
\]
\item Так как $\lambda-\mu\ne0$, делим и получаем $(v,w)=0$. Это определение ортогональности.
\end{enumerate}\finish
\remember{При одинаковых собственных значениях этот довод не даёт ортогональности: внутри одного собственного подпространства ортонормированный базис выбирают отдельно. Для произвольного оператора независимость из раздела 24 есть, а ортогональности может не быть.}

\topic{26}{Критерий Сильвестра}{§ 2.7.2, с. 64--65}
\claim Для вещественной симметричной $n\times n$ матрицы $A$ форма $q(x)=x^TAx$ положительно определена тогда и только тогда, когда все угловые главные миноры $\Delta_1,\ldots,\Delta_n$ положительны.
\idea{Выделить полный квадрат по первой переменной. Оставшаяся форма имеет миноры, выражающиеся через исходные.}
\uses{выделение полного квадрата; блочную формулу определителя; индукцию}
\begin{enumerate}
\item При $n=1$ утверждение очевидно. Пусть оно доказано для размера $n-1$. Запишем
\[
 A=\begin{pmatrix}a&b^T\\b&C\end{pmatrix},\qquad x=(s,z).
\]
При $a>0$ выделение квадрата даёт
\[
 q(s,z)=a\left(s+\frac{b^Tz}{a}\right)^2+z^TBz,
 \quad B=C-\frac{bb^T}{a}.
\]
\item Матрица $B$ симметрична. Для её углового минора размера $j$ блочная формула (или вычитание кратных первой строки) даёт
\[
 \Delta_{j+1}(A)=a\Delta_j(B),\qquad1\le j\le n-1.
\]
\item \textbf{Достаточность.} Если все $\Delta_j(A)>0$, то $a=\Delta_1>0$ и все $\Delta_j(B)>0$. По индукции $z^TBz>0$ при $z\ne0$. В формуле шага 1 при $z\ne0$ положителен второй член, а при $z=0$, $s\ne0$ --- первый. Итак, $q>0$ вне нуля.
\item \textbf{Необходимость.} Если $q$ положительно определена, то $a=q(1,0)>0$. Для $z\ne0$ подставим $s=-b^Tz/a$: получим $z^TBz=q(s,z)>0$. По индукции все миноры $B$ положительны. Формула шага 2 даёт все $\Delta_j(A)>0$.
\end{enumerate}\finish
\remember{Речь именно об угловых главных минорах. Замена строгого знака на $\ge0$ в этом критерии не даёт критерия полуопределённости.}

\topic{27}{Решения по кратному характеристическому корню}{§ 3.2, с. 69--71}
\claim Если $\lambda\in\C$ --- корень кратности $m$ характеристического многочлена $M$ уравнения $M(D)y=0$, $D=d/dx$, то функции
\[
 e^{\lambda x},\ xe^{\lambda x},\ldots,x^{m-1}e^{\lambda x}
\]
являются решениями и линейно независимы.
\idea{Экспонента сдвигает оператор $D$ на $\lambda$, а $m$ производных уничтожают многочлен степени меньше $m$.}
\uses{правило производной произведения; разложение многочлена по корню}
\begin{enumerate}
\item Для любой гладкой $u$
\[
 D(e^{\lambda x}u)=e^{\lambda x}(D+\lambda)u.
\]
Повторяя равенство и складывая с коэффициентами $M$, получаем
\[
 M(D)(e^{\lambda x}u)=e^{\lambda x}M(D+\lambda)u.
\]
\item Кратность означает $M(z)=(z-\lambda)^mQ(z)$. Следовательно,
\[
 M(D+\lambda)=D^mQ(D+\lambda)=Q(D+\lambda)D^m.
\]
Порядок множителей можно менять: оба --- многочлены от одного оператора $D$ с постоянными коэффициентами.
\item Для $u=x^k$, $0\le k<m$, имеем $D^mu=0$. Значит, $M(D)(e^{\lambda x}x^k)=0$.
\item Если $\sum_{k=0}^{m-1}c_kx^ke^{\lambda x}=0$ на интервале, делим на ненулевую экспоненту. Многочлен $\sum c_kx^k$ тождественно нулевой, поэтому все $c_k=0$.
\end{enumerate}\finish
\remember{Доказан блок решений для одного корня. Для вещественного уравнения сопряжённые комплексные корни дают вещественные решения через косинус и синус.}

\topic{28}{Решение системы через матричную экспоненту}{§ 3.3.2, с. 74--76}
\claim Для постоянной вещественной матрицы $A$ единственное решение задачи $x'=Ax$, $x(t_0)=x_0$ имеет вид
\[
 x(t)=e^{(t-t_0)A}x_0,\qquad e^{tA}=\sum_{k=0}^\infty\frac{t^kA^k}{k!}.
\]
\idea{Ряд экспоненты можно дифференцировать; для единственности умножаем любое решение на обратную экспоненту.}
\uses{признак Вейерштрасса; почленное дифференцирование равномерно сходящихся рядов производных}
\begin{enumerate}
\item На $|t|\le R$ нормы членов ряда не больше $(R\|A\|)^k/k!$. Ряд этих чисел сходится. Для ряда производных мажоранта $\|A\|\sum_{j=0}^\infty(R\|A\|)^j/j!$ тоже сходится.
\item Поэтому допустимо дифференцирование:
\[
 \frac d{dt}e^{tA}=Ae^{tA}=e^{tA}A,\qquad e^{0A}=I.
\]
Коммутирование с $A$ следует из того, что каждый член ряда есть степень $A$.
\item Производная $e^{-tA}e^{tA}$ равна нулю по правилу произведения. При $t=0$ произведение равно $I$, значит, $e^{-tA}$ обратна к $e^{tA}$.
\item Предложенная функция удовлетворяет уравнению и начальному условию по шагу 2. Если $z$ --- любое другое решение, то
\[
 \frac d{dt}\bigl(e^{-(t-t_0)A}z(t)\bigr)
 =e^{-(t-t_0)A}(-Az+z')=0.
\]
Следовательно, произведение постоянно и равно $x_0$. Умножаем на $e^{(t-t_0)A}$ и получаем ту же формулу.
\end{enumerate}\finish
\remember{Сходимости самого ряда недостаточно для дифференцирования: в шаге 1 отдельно проверяется ряд производных.}

\topic{29}{Формула Лиувилля--Остроградского}{§ 3.1, с. 66--68}
\claim Для решений $y_1,\ldots,y_n$ уравнения
\[
 y^{(n)}+p_1(t)y^{(n-1)}+\cdots+p_n(t)y=0
\]
с непрерывными коэффициентами на интервале $I$ их вронскиан удовлетворяет
\[
 W(t)=W(t_0)\exp\left(-\int_{t_0}^t p_1(s)\,ds\right).
\]
\idea{При дифференцировании определителя все слагаемые, кроме последнего, имеют одинаковые строки.}
\uses{правило производной определителя; линейность по строке}
\begin{enumerate}
\item Обозначим $R_j=(y_1^{(j)},\ldots,y_n^{(j)})$. Тогда $W=\det(R_0,\ldots,R_{n-1})$, где перечислены строки.
\item Производная определителя --- сумма определителей, в каждом из которых дифференцируется одна строка. Если дифференцируется $R_j$, $j<n-1$, вместо неё возникает уже имеющаяся строка $R_{j+1}$. Такой определитель нулевой.
\item Остаётся производная последней строки. По уравнению
\[
 R_n=-p_1R_{n-1}-p_2R_{n-2}-\cdots-p_nR_0.
\]
По линейности все слагаемые снова имеют повторённую строку, кроме $-p_1R_{n-1}$. Следовательно, $W'=-p_1W$.
\item Чтобы не делить на потенциально нулевой $W$, положим $P(t)=\int_{t_0}^t p_1(s)\,ds$. Тогда
\[
 (e^{P(t)}W(t))'=e^{P(t)}(p_1W+W')=0.
\]
Значит, $e^{P(t)}W(t)=W(t_0)$, что и даёт формулу.
\end{enumerate}\finish
\remember{Формула верна и при $W(t_0)=0$. Вронскиан либо всюду нулевой, либо нигде не нулевой на интервале.}

\topic{30}{Необходимое условие Эйлера}{§ 3.4, с. 76--78}
\claim Пусть $F\in C^2$ в окрестности графика $(t,y_*(t),y_*'(t))$, $y_*\in C^2[a,b]$ --- слабый локальный экстремум
\[
 J[y]=\int_a^bF(t,y(t),y'(t))\,dt
\]
с фиксированными $y(a),y(b)$. Тогда на $(a,b)$
\[
 F_y(t,y_*,y_*')-\frac d{dt}F_p(t,y_*,y_*')=0,\qquad p=y'.
\]
\idea{Изменять функцию вдоль произвольной вариации; одномерная производная в точке экстремума равна нулю.}
\uses{Ферма; дифференцирование интеграла; интегрирование по частям; основную лемму вариационного исчисления}
\begin{enumerate}
\item Берём произвольную $\eta\in C^1[a,b]$ с $\eta(a)=\eta(b)=0$. Функции $y_s=y_*+s\eta$ имеют те же концы и при малом $|s|$ близки к $y_*$ в норме $C^1$.
\item Поэтому $j(s)=J[y_s]$ имеет экстремум при $s=0$, и $j'(0)=0$. Гладкость $F$ и компактность отрезка позволяют дифференцировать под интегралом:
\[
 0=\int_a^b(F_y\eta+F_p\eta')\,dt.
\]
Все производные $F$ вычисляются вдоль $y_*$.
\item Интегрируем второй член по частям. Граничный член $[F_p\eta]_a^b$ равен нулю, поэтому
\[
 \int_a^b\left(F_y-\frac d{dt}F_p\right)\eta\,dt=0
\]
для каждой допустимой $\eta$.
\item Основная лемма говорит: если непрерывная $g$ удовлетворяет $\int g\eta=0$ для всех таких вариаций, то $g=0$. Иначе в малой окрестности точки с $g>0$ (или $g<0$) можно взять неотрицательную ненулевую гладкую вариацию с носителем внутри неё; интеграл имел бы строгий знак. Применяем это к скобке.
\end{enumerate}\finish
\remember{Уравнение Эйлера необходимо, но само по себе не гарантирует экстремума.}

\topic{31}{Формула Байеса}{§ 4.1.2--4.1.3, с. 79--80}
\claim Пусть $(B_i)$ --- конечное или счётное разбиение пространства исходов на события с $P(B_i)>0$, а $P(A)>0$. Тогда
\[
 P(B_k\mid A)=\frac{P(A\mid B_k)P(B_k)}{\sum_iP(A\mid B_i)P(B_i)}.
\]
\idea{Один и тот же совместный исход $A\cap B_k$ записать через две условные вероятности, а знаменатель разложить по гипотезам.}
\uses{определение условной вероятности; счётную аддитивность}
\begin{enumerate}
\item По определению
\[
 P(B_k\mid A)=\frac{P(A\cap B_k)}{P(A)}.
\]
Условие $P(A)>0$ нужно для деления.
\item Поскольку $P(B_k)>0$, другое условное выражение даёт
\[
 P(A\cap B_k)=P(A\mid B_k)P(B_k).
\]
\item События $A\cap B_i$ попарно не пересекаются, а их объединение равно $A$, потому что $B_i$ образуют разбиение. По аддитивности
\[
 P(A)=\sum_iP(A\cap B_i)=\sum_iP(A\mid B_i)P(B_i).
\]
Это одновременно вывод формулы полной вероятности.
\item Подставляя числитель из шага 2 и знаменатель из шага 3 в шаг 1, получаем формулу Байеса.
\end{enumerate}\finish
\remember{Независимость событий здесь не предполагается. Формула пересчитывает вероятности гипотез после наблюдения события $A$.}

\topic{32}{Дисперсия суммы независимых величин}{§ 4.2.1, с. 80--81}
\claim Если $X,Y$ независимы и имеют конечные вторые моменты, то
\[
 \Var(X+Y)=\Var X+\Var Y.
\]
\idea{Раскрыть квадрат центрированной суммы; смешанный член исчезнет благодаря независимости.}
\uses{линейность ожидания; для независимых интегрируемых величин ожидание произведения равно произведению ожиданий}
\begin{enumerate}
\item Обозначим $U=X-\E X$, $V=Y-\E Y$. Конечность вторых моментов обеспечивает конечность всех используемых ожиданий; в частности, $\E|UV|<\infty$ по Коши--Буняковскому.
\item Линейность даёт $\E(X+Y)=\E X+\E Y$. Поэтому
\begin{align*}
 \Var(X+Y)&=\E[(U+V)^2]\\
 &=\E U^2+2\E(UV)+\E V^2.
\end{align*}
\item Вычитание констант сохраняет независимость. Следовательно,
\[
 \E(UV)=(\E U)(\E V)=0\cdot0=0.
\]
\item По определению $\E U^2=\Var X$ и $\E V^2=\Var Y$. Подставляя, получаем требуемое равенство.
\end{enumerate}\finish
\remember{Для зависимых величин появляется $2\operatorname{Cov}(X,Y)$. Независимость --- достаточное условие его исчезновения; достаточно также нулевой ковариации.}

\topic{33}{Неравенство Чебышёва}{§ 4.3, с. 82--83}
\claim Если $\Var X<\infty$, то для каждого $\varepsilon>0$
\[
 P(|X-\E X|\ge\varepsilon)\le\frac{\Var X}{\varepsilon^2}.
\]
\idea{На событии большого отклонения квадрат отклонения не меньше $\varepsilon^2$. Его среднее поэтому контролирует вероятность события.}
\uses{монотонность ожидания; ожидание индикатора $\E\mathbf1_A=P(A)$}
\begin{enumerate}
\item Положим $A=\{|X-\E X|\ge\varepsilon\}$. Для каждого исхода верно
\[
 (X-\E X)^2\ge\varepsilon^2\mathbf1_A.
\]
На $A$ это следует из определения события; вне $A$ правая часть нулевая, а левая неотрицательна.
\item Возьмём ожидания и применим монотонность:
\[
 \Var X=\E(X-\E X)^2
 \ge\varepsilon^2\E\mathbf1_A=\varepsilon^2P(A).
\]
\item Деление на положительное $\varepsilon^2$ даёт утверждение.
\end{enumerate}\finish
\remember{Из этой оценки непосредственно получается закон больших чисел: если независимые $X_j$ имеют общее ожидание $a$ и дисперсии не больше $C$, то $\Var(n^{-1}\sum X_j)\le C/n$, а вероятность отклонения среднего от $a$ не меньше чем на $\varepsilon$ не превосходит $C/(n\varepsilon^2)\to0$.}

\topic{34}{Центральная предельная теорема}{§ 4.4, с. 83--84}
\claim Для независимых одинаково распределённых $X_j$ с $\E X_j=a$ и $0<\Var X_j=\sigma^2<\infty$
\[
 Z_n=\frac{\sum_{j=1}^nX_j-na}{\sigma\sqrt n}
 \xrightarrow{\ d\ }N(0,1).
\]
\idea{Характеристическая функция одного малого слагаемого имеет вид $1-t^2/(2n)+o(1/n)$; независимость превращает сумму в $n$-ю степень.}
\uses{теорему о мажорируемой сходимости; теорему непрерывности Леви; характеристическую функцию $N(0,1)$, равную $e^{-t^2/2}$}
\begin{enumerate}
\item Стандартизуем: $Y_j=(X_j-a)/\sigma$. Тогда $\E Y_j=0$, $\E Y_j^2=1$. Обозначим $\phi(u)=\E e^{iuY_1}$.
\item Для вещественного $v$ положим $r(v)=e^{iv}-1-iv+v^2/2$. По Тейлору $r(v)/v^2\to0$ при $v\to0$; кроме того, $|r(v)|\le v^2$ (остаток после линейного члена не больше $v^2/2$).
\item Следовательно, при $u\to0$ величина $r(uY_1)/u^2\to0$ для каждого исхода и по модулю не больше $Y_1^2$. По мажорируемой сходимости её ожидание стремится к нулю. Итак,
\[
 \phi(u)=1-\frac{u^2}2+o(u^2).
\]
\item По независимости для фиксированного $t$
\[
 \phi_{Z_n}(t)=\bigl[\phi(t/\sqrt n)\bigr]^n
 =\left[1-\frac{t^2}{2n}+o(1/n)\right]^n\longrightarrow e^{-t^2/2}.
\]
Последний предел следует из $\log(1+z)=z+o(z)$ около нуля: логарифм степени стремится к $-t^2/2$.
\item Предельная функция --- характеристическая функция стандартного нормального закона, непрерывная в нуле. По теореме Леви получаем сходимость по распределению.
\end{enumerate}\finish
\remember{Не нужен бесконечный ряд по моментам: второго момента достаточно для одного разложения с остатком $o(u^2)$. Дисперсия должна быть строго положительной.}

\topic{35}{Критерий Коши--Римана}{§ 5.1.1--5.1.2, с. 85--86}
\claim Для $f=u+iv$ комплексная дифференцируемость в $z_0=x_0+iy_0$ равносильна вещественной дифференцируемости $u,v$ в $(x_0,y_0)$ и равенствам
\[
 u_x=v_y,\qquad u_y=-v_x.
\]
При этом $f'(z_0)=u_x+iv_x$; все частные производные берутся в данной точке.
\idea{Комплексное умножение --- специальное вещественное линейное отображение; Коши--Риман описывает его матрицу.}
\uses{определения комплексной и вещественной дифференцируемости}
\begin{enumerate}
\item \textbf{Необходимость.} Если $f'(z_0)=\alpha+i\beta$, то для $h=s+it$
\[
 f(z_0+h)-f(z_0)=(\alpha+i\beta)(s+it)+o(|h|).
\]
Вещественная часть линейного члена равна $\alpha s-\beta t$, мнимая --- $\beta s+\alpha t$. Значит, $u,v$ вещественно дифференцируемы и
\[
 u_x=\alpha,\quad u_y=-\beta,\quad v_x=\beta,\quad v_y=\alpha.
\]
Отсюда требуемые равенства.
\item \textbf{Достаточность.} Из вещественной дифференцируемости
\[
 \Delta f=(u_x+iv_x)s+(u_y+iv_y)t+o(|h|).
\]
По Коши--Риману $u_y+iv_y=-v_x+iu_x=i(u_x+iv_x)$. Поэтому
\[
 \Delta f=(u_x+iv_x)(s+it)+o(|h|).
\]
\item При $h\ne0$ делим на $h$: остаток, делённый на $h$, стремится к нулю по модулю. Значит, комплексная производная существует и равна $u_x+iv_x$.
\end{enumerate}\finish
\remember{Равенств между частными производными без вещественной дифференцируемости недостаточно. Её обеспечивает, например, непрерывность частных производных около точки.}

\topic{36}{Ряд Тейлора регулярной функции}{§ 5.2.2, с. 89--91}
\claim Если $f$ регулярна в диске $|z-z_0|<R$, то там
\[
 f(z)=\sum_{n=0}^\infty c_n(z-z_0)^n,\qquad
 c_n=\frac{f^{(n)}(z_0)}{n!}.
\]
Сходимость абсолютная и равномерная на каждом меньшем замкнутом диске.
\idea{В формуле Коши разложить ядро $1/(w-z)$ в геометрический ряд.}
\uses{интегральную формулу Коши и её формулу для производных; геометрический ряд}
\begin{enumerate}
\item Фиксируем $r<\rho<R$ и интегрируем по положительной окружности $\gamma:|w-z_0|=\rho$. При $|z-z_0|\le r$ отношение $|(z-z_0)/(w-z_0)|\le r/\rho<1$.
\item Поэтому равномерно по $w\in\gamma$ и $|z-z_0|\le r$
\[
 \frac1{w-z}=\sum_{n=0}^\infty\frac{(z-z_0)^n}{(w-z_0)^{n+1}}.
\]
\item Подставляем в формулу Коши. Равномерная сходимость позволяет интегрировать почленно:
\[
 f(z)=\sum_{n=0}^\infty(z-z_0)^n
 \frac1{2\pi i}\oint_\gamma\frac{f(w)}{(w-z_0)^{n+1}}\,dw.
\]
По формуле Коши для производных интеграл-коэффициент равен $f^{(n)}(z_0)/n!$.
\item Если $M_\rho=\max_\gamma|f|$, оценка интеграла даёт $|c_n|\le M_\rho/\rho^n$. Следовательно,
\[
 |c_n(z-z_0)^n|\le M_\rho(r/\rho)^n.
\]
Геометрическая мажоранта доказывает абсолютную и равномерную сходимость. Любую точку исходного диска можно включить в такой меньший диск.
\end{enumerate}\finish
\remember{Окружность интегрирования выбирается строго внутри области регулярности, но снаружи всех точек, для которых нужна общая оценка.}

\topic{37}{Единственность коэффициентов Лорана}{§ 5.3.1, с. 91--92}
\claim Если в кольце $r<|z-a|<R$ функция представлена рядом Лорана
\[
 f(z)=\sum_{n=-\infty}^{\infty}c_n(z-a)^n
\]
с локально равномерной сходимостью, то для любого $r<\rho<R$ и любого целого $k$
\[
 c_k=\frac1{2\pi i}\oint_{|z-a|=\rho}\frac{f(z)}{(z-a)^{k+1}}\,dz.
\]
В частности, коэффициенты определяются функцией единственным образом. Окружность ориентирована против часовой стрелки.
\idea{Интегрирование уничтожает все степени, кроме $(z-a)^{-1}$.}
\uses{почленное интегрирование равномерно сходящихся рядов; параметризацию окружности}
\begin{enumerate}
\item На выбранной окружности оба ряда --- по неотрицательным и по отрицательным степеням --- сходятся равномерно. Умножение на непрерывную ограниченную функцию $(z-a)^{-k-1}$ сохраняет равномерную сходимость. Поэтому интегрируем почленно.
\item Для любого целого $m$ при $z=a+\rho e^{it}$, $0\le t\le2\pi$, имеем
\[
 \oint(z-a)^m\,dz=i\rho^{m+1}\int_0^{2\pi}e^{i(m+1)t}\,dt
 =\begin{cases}2\pi i,&m=-1,\\0,&m\ne-1.\end{cases}
\]
\item Следовательно, в интеграле от ряда остаётся только член с $n-k-1=-1$, то есть $n=k$. Значение интеграла равно $2\pi i\,c_k$, что даёт формулу.
\item Правая часть зависит лишь от $f$. Два ряда Лорана для одной функции должны иметь одинаковые коэффициенты.
\end{enumerate}\finish
\remember{Доказывается формула коэффициентов и единственность при существующем разложении; существование ряда Лорана здесь не выводится.}

\topic{38}{Теорема о вычетах}{§ 5.4.1--5.4.2, с. 94--96}
\claim Пусть $D$ --- ограниченная область с простой замкнутой кусочно-гладкой границей, а $f$ регулярна в окрестности $\overline D$, кроме конечного числа изолированных особых точек $a_1,\ldots,a_m\in D$. Тогда при положительном обходе
\[
 \oint_{\partial D}f(z)\,dz=2\pi i\sum_{j=1}^m\operatorname{res}_{a_j}f.
\]
\idea{Вырезать маленькие диски вокруг особенностей. В оставшейся области применить Коши; интегралы по маленьким окружностям выделят вычеты.}
\uses{интегральную теорему Коши для полной границы области с отверстиями; ряд Лорана; раздел 37}
\begin{enumerate}
\item Выбираем попарно непересекающиеся замкнутые диски $\overline B(a_j,r_j)\subset D$, каждый из которых содержит только одну особую точку. Это возможно, поскольку точек конечное число и все они внутренние.
\item Обозначим их окружности через $\gamma_j$, ориентированные против часовой стрелки, и вырежем диски из $D$. На окрестности замыкания оставшейся области $f$ регулярна.
\item Её положительная граница состоит из внешней $\partial D$ и окружностей отверстий, обходимых по часовой стрелке. Поэтому теорема Коши даёт
\[
 0=\oint_{\partial D}f(z)\,dz-\sum_{j=1}^m\oint_{\gamma_j}f(z)\,dz.
\]
\item Вокруг $a_j$ есть ряд Лорана. Вычет --- его коэффициент $c_{-1}$. По формуле коэффициентов (раздел 37, $k=-1$)
\[
 \oint_{\gamma_j}f(z)\,dz=2\pi i\,c_{-1}
 =2\pi i\operatorname{res}_{a_j}f.
\]
\item Подставляем в шаг 3 и переносим сумму вправо. Получаем утверждение.
\end{enumerate}\finish
\remember{Минус перед внутренними окружностями возникает из ориентации отверстий. Особенностей на исходном контуре быть не должно.}

\end{document}

